\documentclass[11pt]{article}
\usepackage[margin=1in]{geometry}
\usepackage{amsmath,amssymb,amsthm,mathtools}
\usepackage{booktabs}
\usepackage{hyperref}
\usepackage{algorithm,algorithmic}
\usepackage{bm}
\usepackage{xcolor}
\usepackage{tcolorbox}

\newtheorem{theorem}{Theorem}[section]
\newtheorem{proposition}[theorem]{Proposition}
\newtheorem{lemma}[theorem]{Lemma}
\newtheorem{corollary}[theorem]{Corollary}
\newtheorem{definition}[theorem]{Definition}
\newtheorem{assumption}[theorem]{Assumption}
\newtheorem{remark}[theorem]{Remark}

\newcommand{\R}{\mathbb{R}}
\newcommand{\E}{\mathbb{E}}
\newcommand{\Var}{\mathrm{Var}}
\newcommand{\Cov}{\mathrm{Cov}}
\newcommand{\tr}{\mathrm{tr}}
\newcommand{\diag}{\mathrm{diag}}
\newcommand{\sign}{\mathrm{sign}}
\newcommand{\MAD}{\mathrm{MAD}}
\newcommand{\EWMA}{\mathrm{EWMA}}
\newcommand{\clip}{\mathrm{clip}}
\newcommand{\bw}{\bm{w}}
\newcommand{\bmu}{\bm{\mu}}
\newcommand{\bkappa}{\bm{\kappa}}

\title{\textbf{PRISM: Mathematical Foundations}\\[6pt]
\large Sharpe--Optimal Allocation under Forecast Uncertainty\\
with Kantorovich Order Flow}
\author{Benjamin Coriat\\McGill University --- ECSE 551}
\date{April 2026}

\begin{document}
\maketitle

\begin{abstract}
We derive from first principles the optimal portfolio allocation rule when the investor
possesses \emph{noisy} return and covariance forecasts together with a full characterisation
of the forecast errors.  The key idea is that the classical mean--variance programme
is mis-specified when the inputs are themselves uncertain: we show how to build a
\emph{confidence-weighted} expected return vector and a \emph{Bayesian-blended}
covariance matrix, then prove that the resulting Sharpe-maximising programme is a
second-order cone programme.  Finally, we derive the Kantorovich optimal transport
plan that converts the current portfolio into the target at minimum transaction cost,
and show that its entropically regularised form admits a closed-form Sinkhorn iteration.
\end{abstract}

\newpage
\tableofcontents
\newpage

%======================================================================
\section{Setup and Data}\label{sec:setup}
%======================================================================

Fix a universe of $n$ assets and a weekly time index $t$ (Friday close).

\begin{definition}[Returns and volatility]\label{def:basics}
For asset $i$ at week $t$:
\begin{align}
r_{i,t} &:= \frac{P_{i,t}}{P_{i,t-1}} - 1 && \text{(simple weekly return)},\\[4pt]
\hat\sigma_{i,t}^{\mathrm{GK}} &:= \sqrt{\max\!\Big(\tfrac12(\ln H - \ln L)^2
  - (2\ln 2-1)(\ln C - \ln O)^2,\;0\Big)} && \text{(Garman--Klass vol)}.
\end{align}
\end{definition}

\begin{remark}[Why Garman--Klass?]
Close-to-close volatility uses only two prices per period.  The GK estimator
uses all four OHLC prices; under driftless GBM one can show
$\Var(\hat\sigma_{\mathrm{GK}}^2)\approx 0.273\,\sigma^4$,
roughly $7.3\times$ more efficient.  Better volatility estimates
feed directly into a better covariance matrix, so this matters for allocation.
\end{remark}

\begin{definition}[Rolling correlation]
With window $W=26$,
\[
\hat\rho_{ij,t} := \frac{\sum_{s=t-W+1}^{t}(r_{i,s}-\bar r_i)(r_{j,s}-\bar r_j)}
  {\sqrt{\sum_s(r_{i,s}-\bar r_i)^2\sum_s(r_{j,s}-\bar r_j)^2}}.
\]
Collect these into the correlation matrix $\hat C_t\in\R^{n\times n}$.
\end{definition}

%======================================================================
\section{The Forecasting Step (Black Box)}\label{sec:forecast}
%======================================================================

A stacked ensemble regressor (Ridge, GBM, Random Forest, MLP, etc.)
produces, at each date $t$:
\begin{enumerate}
\item A predicted return vector $\hat\bmu_{t+1}\in\R^n$.
\item A predicted correlation matrix $\hat C_{t+1}\in\R^{n\times n}$.
\end{enumerate}
The details of the ensemble are not the focus here.  What matters is that
these forecasts are \emph{imperfect}, and we have a way to measure how imperfect
they are --- the \textbf{error characterisation} described next.

%======================================================================
\section{Error Characterisation}\label{sec:error}
%======================================================================

For each asset $i$, define the signed forecast error at past date $s$:
\[
\varepsilon_{i,s} := \hat r_{i,s}^{\mathrm{pred}} - r_{i,s+1}.
\]
We maintain a running history $\mathcal{H}_{i,t}=\{\varepsilon_{i,s}\}_{s<t}$
and compute five statistics that together describe \emph{how trustworthy}
the forecast for asset~$i$ is.

\begin{definition}[Return error matrix $E_\mu\in\R^{n\times 5}$]\label{def:Emu}
\begin{align}
[E_\mu]_{i,1} &:= \sum_s w_s\,\varepsilon_{i,s}
  &&\text{EWMA bias ($\lambda=0.94$)},\label{eq:bias}\\
[E_\mu]_{i,2} &:= \mathbf{1}[\sign(\hat r_{i,t-1})=\sign(r_{i,t})]
  &&\text{last directional hit},\label{eq:dir}\\
[E_\mu]_{i,3} &:= \mathrm{std}(\mathcal{H}_{i,t})
  &&\text{error volatility},\label{eq:evol}\\
[E_\mu]_{i,4} &:= \mathrm{skew}(\mathcal{H}_{i,t})
  &&\text{error skewness},\label{eq:eskew}\\
[E_\mu]_{i,5} &:= \mathrm{exkurt}(\mathcal{H}_{i,t})
  &&\text{error excess kurtosis}.\label{eq:ekurt}
\end{align}
\end{definition}

\begin{remark}[Intuition]
Column~1 tells us whether the model is \emph{systematically} over- or under-predicting
(and recent errors are up-weighted).
Column~2 says whether it at least gets the \emph{sign} right.
Columns~3--5 describe the \emph{shape} of the error distribution: a wide, left-skewed,
heavy-tailed error distribution means we should trust the forecast less.
\end{remark}

Similarly, for each pair $(i,j)$ we track $\varepsilon_{ij,s}^C:=\hat\rho_{ij,s}^{\mathrm{pred}}-\hat\rho_{ij,s+1}$
and store bias, std, skewness, kurtosis in a tensor $E_C\in\R^{n\times n\times 4}$.

Finally, the last $k=26$ weeks of GK volatilities are stored as $V\in\R^{n\times k}$.

%======================================================================
\section{From Error Statistics to Confidence}\label{sec:confidence}
%======================================================================

The central question: given $E_\mu$, how should we \emph{down-weight}
the return forecasts we don't trust?

\subsection{Designing the confidence score}

We want a score $\kappa_i\in(0,1]$ for each asset, where $\kappa_i=1$ means
``fully trust the forecast'' and $\kappa_i\approx 0$ means ``ignore it.''
We build $\kappa_i$ multiplicatively from four penalty terms, each in $(0,1]$.

\begin{tcolorbox}[colback=blue!3, colframe=blue!40!black, title=Design principle]
Each penalty is an exponential decay in a normalised badness measure.
If an error statistic is at the \emph{median} level, the penalty is moderate;
if it is many MADs above median, the penalty drives $\kappa_i$ toward zero.
\end{tcolorbox}

\begin{definition}[Component scores]\label{def:scores}
Let $b_i:=[E_\mu]_{i,1}$, $d_i:=[E_\mu]_{i,2}$, $v_i:=[E_\mu]_{i,3}$,
$\xi_i:=[E_\mu]_{i,4}$, $\kappa_i^{\mathrm{ex}}:=\max([E_\mu]_{i,5},0)$.

\medskip\noindent
\textbf{1.\ Bias penalty.}  Large systematic bias $\Rightarrow$ low trust:
\begin{equation}\label{eq:sbias}
s_i^{\mathrm{bias}} := \exp\!\bigg(-\alpha_1\,\frac{|b_i|}{\MAD(\bm{b})}\bigg),
\qquad \MAD(\bm{b}):=\max\!\big(\mathrm{median}(|b_j-\mathrm{median}(\bm{b})|),\,10^{-8}\big).
\end{equation}
\emph{Why MAD?}  The median absolute deviation is robust to outliers,
so one rogue asset doesn't inflate the scale and mask problems in others.

\medskip\noindent
\textbf{2.\ Directional accuracy.}  If the model can't even get the sign right, down-weight:
\begin{equation}\label{eq:sdir}
s_i^{\mathrm{dir}} := \beta_1\, d_i + (1-\beta_1),
\qquad \beta_1=0.3.
\end{equation}
This is simply a linear interpolation: a correct sign gives $s^{\mathrm{dir}}=1$,
a miss gives $1-\beta_1=0.7$.  The penalty is intentionally mild because a single
wrong sign shouldn't nuke the confidence.

\medskip\noindent
\textbf{3.\ Dispersion penalty.}  High error volatility, especially combined with fat tails,
means the forecast is unreliable:
\begin{equation}\label{eq:sdisp}
\delta_i := v_i\,(1+\gamma_1\kappa_i^{\mathrm{ex}}),\qquad
s_i^{\mathrm{disp}} := \exp\!\bigg(-\alpha_2\,\frac{\delta_i}{\mathrm{median}(\bm\delta)}\bigg).
\end{equation}
The term $\gamma_1\kappa_i^{\mathrm{ex}}$ inflates the effective dispersion when the error
distribution has excess kurtosis --- because fat tails mean the standard deviation
\emph{underestimates} the true forecast risk.

\medskip\noindent
\textbf{4.\ Skewness penalty.}  Negatively skewed errors (model over-predicts) are
penalised:
\begin{equation}\label{eq:sskew}
s_i^{\mathrm{skew}} := \exp\!\big(-\alpha_3\max(-\xi_i,0)\big).
\end{equation}
Only \emph{negative} skew is penalised: over-prediction is worse than under-prediction
for a long-biased portfolio.
\end{definition}

\begin{definition}[Composite confidence]\label{def:kappa}
\begin{equation}\label{eq:kappa}
\kappa_i := \clip\!\bigg(
  \prod_{k=1}^{4}\big(s_i^{(k)}\big)^{\omega_k},\;10^{-4},\;1\bigg),
\qquad \omega_k=\tfrac14.
\end{equation}
\end{definition}

\begin{remark}[Why geometric mean?]
A multiplicative structure means that \emph{any single} bad signal can substantially
reduce confidence.  If we used an arithmetic mean, a terrible bias score could be
offset by a good directional score, which is undesirable: we want a conjunction
(``trust only if \emph{all} diagnostics are acceptable'').
\end{remark}

\subsection{Confidence-weighted returns}

Armed with $\kappa_i$, we shrink each forecast toward the cross-sectional mean:
\begin{equation}\label{eq:mu_tilde}
\tilde\mu_i := \kappa_i\,\hat\mu_i + (1-\kappa_i)\,\mu_0,
\qquad \mu_0:=\frac{1}{n}\sum_{j=1}^n \hat\mu_j.
\end{equation}

\begin{remark}[Interpretation]
When $\kappa_i=1$, we fully believe the forecast.
When $\kappa_i\to 0$, the asset's expected return is pulled to the market average,
reflecting maximum uncertainty (``we know nothing special about this asset'').
This is a form of \emph{James--Stein shrinkage}: it provably reduces expected
squared error relative to the raw forecasts when $n\ge 3$.
\end{remark}

%======================================================================
\section{Bayesian Posterior Estimation}\label{sec:bayes}
%======================================================================

The confidence-weighted returns $\tilde\bmu$ are a point estimate.
We now embed them in a Bayesian framework to produce a \emph{regularised}
mean and covariance that account for estimation uncertainty.

\subsection{Prior specification}

\begin{assumption}[Normal--Inverse-Wishart conjugate prior]\label{ass:niw}
\begin{align}
\Sigma &\sim \mathcal{IW}(\nu_0,\Psi_0), &
  \nu_0 &= n+2,\\
\bmu\mid\Sigma &\sim \mathcal{N}(\bm{m}_0,\,\tau_0^{-1}\Sigma), &
  \tau_0 &= 1,\quad \bm{m}_0=\mu_0\bm{1}.
\end{align}
The scale matrix is $\Psi_0=\max(\nu_0-n-1,1)\,\hat\Sigma_{\mathrm{hist}}$,
where $\hat\Sigma_{\mathrm{hist}}$ is the historical sample covariance
on the training period.
\end{assumption}

\begin{remark}[Why these choices?]
Setting $\nu_0=n+2$ gives the Inverse-Wishart the minimum number of
degrees of freedom for which the mean exists --- i.e.\ a \emph{diffuse}
prior that doesn't impose strong beliefs.
$\tau_0=1$ means the prior on $\bmu$ has the same weight as a single
observation, so the data quickly dominates.
$\bm{m}_0=\mu_0\bm{1}$ centres the prior on equal expected returns
(maximum ignorance).
\end{remark}

\subsection{Posterior update}

We treat the confidence-weighted vector $\tilde\bmu$ as a single
``observation'' from $\mathcal{N}(\bmu,\Sigma)$, but re-weighted by the
confidence matrix $K=\diag(\kappa_1,\dots,\kappa_n)$.

\begin{proposition}[Posterior point estimates]\label{prop:posterior}
The posterior mean and covariance mode are:
\begin{align}
\bmu_{t+1}^B &= \frac{\tau_0\,\bm{m}_0 + K\tilde\bmu}{\tau_0+1},\label{eq:muB}\\[6pt]
\Sigma_{t+1}^B &= \frac{\Psi_n}{\max(\nu_0+1-n-1,\,1)},\label{eq:SigB}
\end{align}
where $\Psi_n:=\Psi_0+\frac{\tau_0}{\tau_0+1}(\tilde\bmu-\bm{m}_0)(\tilde\bmu-\bm{m}_0)^\top$.
\end{proposition}

\begin{proof}[Derivation]
Under the NIW conjugate model with prior $(\bm{m}_0,\tau_0,\Psi_0,\nu_0)$,
the standard update after observing data $\bar{\bm{x}}$ with effective
sample size $n_{\mathrm{eff}}$ gives posterior parameters
$\bm{m}_n = \frac{\tau_0 \bm{m}_0 + n_{\mathrm{eff}}\bar{\bm{x}}}{\tau_0+n_{\mathrm{eff}}}$.
Here we have a single ``observation'' ($n_{\mathrm{eff}}=1$) but it is
modulated by the confidence matrix $K$: assets with $\kappa_i\approx 1$
contribute their forecast at full strength, while low-confidence assets
are effectively shrunk toward the prior mean $\mu_0$.

The scale matrix update
$\Psi_n = \Psi_0 + \frac{\tau_0}{\tau_0+1}(\tilde\bmu-\bm{m}_0)(\tilde\bmu-\bm{m}_0)^\top$
is the standard NIW scatter update: it inflates the posterior covariance
in directions where the observation deviates from the prior mean,
reflecting the additional uncertainty from disagreement between
prior and data.

The result is projected onto the positive-definite cone via
Higham alternating projections (nearest correlation/covariance matrix).
\end{proof}

%======================================================================
\section{Robust Covariance Construction}\label{sec:cov}
%======================================================================

The Bayesian posterior $\Sigma_{t+1}^B$ handles mean uncertainty.
We also need a \emph{forward-looking} covariance from the predicted
correlation matrix and recent volatilities.

\subsection{EWMA volatility}

From the volatility history $V\in\R^{n\times k}$ ($k=26$), with decay $\lambda_v=0.94$:
\begin{equation}\label{eq:ewma_vol}
\hat\sigma_i := \sqrt{\sum_{\ell=1}^k w_\ell\,V_{i,\ell}^2},
\qquad w_\ell = \frac{(1-\lambda_v)\lambda_v^{k-\ell}}{1-\lambda_v^k}.
\end{equation}

\begin{remark}
The EWMA weights recent observations exponentially more.
$\lambda_v=0.94$ gives a half-life of about $\ln 2/\ln(1/0.94)\approx 11$ weeks,
so the estimate adapts to regime changes within a quarter.
\end{remark}

\subsection{Regime adjustment}

If current vol is elevated relative to its own $k$-week average,
we further inflate it:
\begin{equation}\label{eq:regime}
\hat\sigma_i^{\mathrm{adj}} := \hat\sigma_i\cdot\big(1+\eta\max(r_i^{\mathrm{vol}}-1,0)\big),
\qquad r_i^{\mathrm{vol}}:=\hat\sigma_i/\bar V_i,\quad \eta=0.5.
\end{equation}
This is a simple asymmetric adjustment: vol is only inflated (never deflated),
which biases toward conservatism during stress.

\subsection{Final blended covariance}

The forward-looking covariance uses the predicted correlation
$\hat C_{t+1}$ and adjusted vols:
\begin{equation}
\Sigma^{\mathrm{fwd}} := D_\sigma\,\hat C_{t+1}\,D_\sigma,
\qquad D_\sigma=\diag(\hat\sigma_1^{\mathrm{adj}},\dots,\hat\sigma_n^{\mathrm{adj}}).
\end{equation}

We blend with the Bayesian posterior:
\begin{equation}\label{eq:Sigma_final}
\boxed{\Sigma_{\mathrm{final}} := \zeta\,\Sigma_{t+1}^B + (1-\zeta)\,\Sigma^{\mathrm{fwd}},
\qquad \zeta=0.5.}
\end{equation}

This is a convex combination of two PSD matrices, hence PSD.
We project to strictly PD (eigenvalue floor) for numerical stability.

%======================================================================
\section{Sharpe-Optimal Allocation}\label{sec:sharpe}
%======================================================================

\subsection{The classical Markowitz problem and its flaw}

The textbook approach maximises the Sharpe ratio:
\begin{equation}\label{eq:SR}
\mathrm{SR}(\bw) := \frac{\bw^\top\bmu-r_f}{\sqrt{\bw^\top\Sigma\bw}}.
\end{equation}

The well-known difficulty is that the solution $\bw^*\propto\Sigma^{-1}\bmu$
is \emph{extremely sensitive} to errors in $\bmu$ and $\Sigma$.
Small perturbations in the inputs produce wildly different portfolios
(Michaud, 1989).

\begin{tcolorbox}[colback=green!3, colframe=green!40!black, title=Our remedy]
By replacing $\bmu$ with the confidence-shrunk, Bayesian-updated
$\bmu_{t+1}^B$ and $\Sigma$ with the blended $\Sigma_{\mathrm{final}}$,
we have already regularised the inputs.  The optimiser now sees
``de-noised'' inputs whose estimation error has been attenuated.
This is the payoff of Sections~\ref{sec:confidence}--\ref{sec:cov}.
\end{tcolorbox}

\subsection{The constrained programme}

We solve, for each risk-aversion $\gamma>0$:
\begin{equation}\label{eq:qp}
\bw^*(\gamma) := \arg\max_{\bw\in\R^n}\;
  \underbrace{(\bmu_{t+1}^B - r_f\bm{1})^\top\bw}_{\text{expected excess return}}
  -\;\gamma\underbrace{\bw^\top\Sigma_{\mathrm{final}}\bw}_{\text{portfolio variance}}
\end{equation}
subject to
\begin{align}
\bm{1}^\top\bw &= 1 && \text{(fully invested)},\label{eq:budget}\\
\|\bw\|_1 &\le L_{\max}=1.6 && \text{(gross leverage cap)},\label{eq:lev}\\
|w_i| &\le w_{\max}=0.20,\;\forall\,i && \text{(concentration limit)}.\label{eq:conc}
\end{align}

\begin{remark}[Why not maximise SR directly?]
Maximising $\mathrm{SR}(\bw)$ is a fractional programme.
The standard trick (Cornuejols \& T\"ut\"unc\"u) is to fix the denominator
and sweep $\gamma$.  For each $\gamma$, \eqref{eq:qp} is a convex QP
(since $\Sigma_{\mathrm{final}}\succ 0$), solvable in milliseconds by SCS.
We then pick the $\gamma^*$ that achieves the highest ex-ante Sharpe:
\begin{equation}
\gamma^* := \arg\max_{\gamma\in\Gamma}\;\mathrm{SR}(\bw^*(\gamma)),
\qquad \bw_{t+1}^* := \bw^*(\gamma^*).
\end{equation}
In practice we use 50 log-spaced values $\gamma\in[10^{-2},10^3]$.
\end{remark}

\subsection{Deriving the KKT conditions}

The Lagrangian of \eqref{eq:qp} (ignoring the $\ell_1$ and box constraints
for clarity) is
\begin{equation}
\mathcal{L}(\bw,\lambda) = (\bmu_{t+1}^B-r_f\bm{1})^\top\bw
  -\gamma\bw^\top\Sigma_{\mathrm{final}}\bw - \lambda(\bm{1}^\top\bw-1).
\end{equation}
Setting $\nabla_{\bw}\mathcal{L}=0$:
\begin{equation}
\bmu_{t+1}^B - r_f\bm{1} - 2\gamma\Sigma_{\mathrm{final}}\bw - \lambda\bm{1} = 0
\quad\Longrightarrow\quad
\bw = \frac{1}{2\gamma}\Sigma_{\mathrm{final}}^{-1}(\bmu_{t+1}^B - r_f\bm{1} - \lambda\bm{1}).
\end{equation}
Substituting into $\bm{1}^\top\bw=1$ and solving for $\lambda$:
\begin{equation}
\lambda = \frac{\bm{1}^\top\Sigma_{\mathrm{final}}^{-1}(\bmu_{t+1}^B-r_f\bm{1})
  - 2\gamma}{\bm{1}^\top\Sigma_{\mathrm{final}}^{-1}\bm{1}}.
\end{equation}

This is the classical two-fund separation result.  The additional
$\ell_1$ and box constraints are handled by the QP solver (they add
inequality constraints and corresponding dual variables, but the
problem remains convex).

%======================================================================
\section{Volatility Targeting}\label{sec:voltarget}
%======================================================================

After allocation, we scale returns to a target annualised volatility
$\sigma^*=0.15$ (15\%):
\begin{equation}\label{eq:vt}
r_t^{\mathrm{VT}} := r_t^{\mathrm{raw}}\cdot\ell_t,
\qquad \ell_t:=\min\!\bigg(\frac{\sigma^*}
  {\hat\sigma_t^{\mathrm{port}}},\;\ell_{\max}\bigg),
\end{equation}
where $\hat\sigma_t^{\mathrm{port}}=\mathrm{std}(\{r_s^{\mathrm{raw}}\}_{s=t-25}^t)\cdot\sqrt{52}$
and $\ell_{\max}=2.5$.

\begin{remark}
This ensures the backtest is comparable across regimes.
During calm periods, leverage is increased (up to $2.5\times$);
during crises, exposure is cut.  The 26-week lookback means the
adjustment is smooth and avoids whipsawing.
\end{remark}

%======================================================================
\section{Optimal Transport for Order Execution}\label{sec:ot}
%======================================================================

We now have a target portfolio $\bw_{t+1}^*$ and a current portfolio $\bw_t$.
The rebalancing trade $\Delta\bw=\bw_{t+1}^*-\bw_t$ incurs transaction costs.
The question is: \emph{what is the cheapest way to rearrange the current portfolio
into the target?}

This is precisely a \textbf{mass transportation problem}.

\subsection{The Kantorovich formulation}

Think of the current weights $\bw_t$ as a ``pile of sand'' distributed
across $n$ assets, and the target $\bw_{t+1}^*$ as the desired landscape.
A \emph{transport plan} $T\in\R^{n\times n}_{\ge 0}$ specifies how much
weight to move from asset $i$ to asset $j$: $T_{ij}$ is the mass
shipped from $i$ to $j$.

\begin{definition}[Transport polytope]
The set of valid plans:
\begin{equation}
\Pi(\bm{a},\bm{b}) := \{T\ge 0 : T\bm{1}=\bm{a},\;T^\top\bm{1}=\bm{b}\},
\end{equation}
where $\bm{a}:=|\bw_t|/\||\bw_t|\|_1$ and $\bm{b}:=|\bw_{t+1}^*|/\||\bw_{t+1}^*|\|_1$
are the normalised absolute weights (probability distributions on assets).
\end{definition}

\begin{definition}[Cost matrix]
The cost of shipping one unit from $i$ to $j$:
\begin{equation}\label{eq:cost}
C_{ij} := \begin{cases}0 & i=j,\\ |c_i|+|c_j| & i\neq j,\end{cases}
\end{equation}
where $c_i$ is the proportional transaction cost for asset $i$
(default: 10\,bps).  The $i=j$ case is free because holding an asset
costs nothing; the $i\neq j$ case is the cost of selling $i$ and buying $j$.
\end{definition}

\begin{definition}[Kantorovich optimal transport]\label{def:kant}
\begin{equation}\label{eq:kant}
T^* := \arg\min_{T\in\Pi(\bm{a},\bm{b})}\;\langle C,T\rangle
  = \arg\min_{T\in\Pi(\bm{a},\bm{b})}\sum_{i,j}C_{ij}T_{ij}.
\end{equation}
\end{definition}

This is a linear programme and can be solved exactly.  However,
for larger $n$ or frequent rebalancing, we prefer the faster
\emph{entropically regularised} version.

\subsection{Entropic regularisation}

\begin{proposition}[Regularised problem]\label{prop:entropic}
Adding an entropic penalty:
\begin{equation}\label{eq:ot_reg}
T_\varepsilon^* := \arg\min_{T\in\Pi(\bm{a},\bm{b})}
  \langle C,T\rangle + \varepsilon\sum_{i,j}T_{ij}(\log T_{ij}-1),
\end{equation}
with $\varepsilon = 0.05\cdot\mathrm{median}(\{C_{ij}:i\neq j,\,C_{ij}>0\})$.
\end{proposition}

\begin{proof}[Derivation of the solution structure]
Write the Lagrangian with dual variables $\bm{f}\in\R^n$, $\bm{g}\in\R^n$
for the marginal constraints:
\begin{equation}
\mathcal{L}(T,\bm{f},\bm{g}) = \sum_{ij}C_{ij}T_{ij}
  + \varepsilon\sum_{ij}T_{ij}(\log T_{ij}-1)
  + \sum_i f_i\Big(a_i - \sum_j T_{ij}\Big)
  + \sum_j g_j\Big(b_j - \sum_i T_{ij}\Big).
\end{equation}
Setting $\partial\mathcal{L}/\partial T_{ij}=0$:
\begin{equation}
C_{ij} + \varepsilon\log T_{ij} - f_i - g_j = 0
\quad\Longrightarrow\quad
T_{ij} = \exp\!\Big(\frac{f_i-C_{ij}+g_j}{\varepsilon}\Big).
\end{equation}
Define $u_i:=e^{f_i/\varepsilon}$, $v_j:=e^{g_j/\varepsilon}$,
$K_{ij}:=e^{-C_{ij}/\varepsilon}$ (the \emph{Gibbs kernel}).  Then:
\begin{equation}\label{eq:Tstar_form}
\boxed{T_{ij}^* = u_i\,K_{ij}\,v_j.}
\end{equation}
The vectors $\bm{u},\bm{v}$ are determined by the marginal constraints
$T\bm{1}=\bm{a}$ and $T^\top\bm{1}=\bm{b}$, which become:
\begin{equation}
u_i\sum_j K_{ij}v_j = a_i,\qquad
v_j\sum_i K_{ij}u_i = b_j.
\end{equation}
These are solved by alternating normalisation --- the Sinkhorn algorithm.
\end{proof}

\subsection{Sinkhorn algorithm}

\begin{algorithm}[h]
\caption{Sinkhorn iterations}
\begin{algorithmic}[1]
\STATE \textbf{Input:} marginals $\bm{a},\bm{b}\in\Delta^{n-1}$, kernel $K$,
  tolerance $\delta=10^{-8}$
\STATE $\bm{v}\leftarrow\bm{1}$
\FOR{$\ell=1,\dots,200$}
  \STATE $\bm{u}\leftarrow \bm{a}\oslash(K\bm{v})$
  \hfill\textit{// enforce row marginal}
  \STATE $\bm{v}\leftarrow \bm{b}\oslash(K^\top\bm{u})$
  \hfill\textit{// enforce column marginal}
  \IF{$\|K^\top\bm{u}\odot\bm{v}-\bm{b}\|_\infty<\delta$}
    \STATE \textbf{break}
  \ENDIF
\ENDFOR
\RETURN $T^*=\diag(\bm{u})\,K\,\diag(\bm{v})$
\end{algorithmic}
\end{algorithm}

\begin{proposition}[Convergence]
Each Sinkhorn iteration is a coordinate ascent step on the dual problem.
The dual objective is strictly concave (since $\varepsilon>0$), so the iterates
converge linearly to the unique optimum at rate $O(e^{-c/\varepsilon})$.
In practice, convergence occurs in 10--30 iterations.
\end{proposition}

\subsection{Transaction cost and net return}

The optimal transport cost is
\begin{equation}
\mathcal{C}^* = \sum_{i\neq j}C_{ij}\,T_{ij}^*,
\end{equation}
and the net portfolio return is
\begin{equation}
r_t^{\mathrm{raw}} = (\bw_{t+1}^*)^\top\bm{r}_t - \mathcal{C}^*.
\end{equation}

%======================================================================
\section{Performance Metrics}\label{sec:metrics}
%======================================================================

Let $\{r_t\}_{t=1}^T$ be the final (vol-targeted) return series with
equity curve $\Pi_t=\prod_{s=1}^t(1+r_s)$.

\begin{align}
\text{Total return} &:= \Pi_T - 1,\\
\text{Annualised return} &:= \Pi_T^{52/T} - 1,\\
\text{Annualised vol} &:= \mathrm{std}(\{r_t\})\cdot\sqrt{52},\\
\text{Sharpe ratio} &:= \frac{\text{Ann.\ return}}{\text{Ann.\ vol}},\\
\text{Max drawdown} &:= \min_t\;\frac{\Pi_t - \max_{s\le t}\Pi_s}{\max_{s\le t}\Pi_s},\\
\text{Calmar ratio} &:= \frac{\text{Ann.\ return}}{|\text{MDD}|},\\
\text{Sortino ratio} &:= \frac{\text{Ann.\ return}}
  {\mathrm{std}(\{r_t:r_t<0\})\cdot\sqrt{52}}.
\end{align}

\newpage
%======================================================================
\section*{References}
%======================================================================

\begin{enumerate}
\item Garman, M.B.\ \& Klass, M.J.\ (1980). On the estimation of security price volatilities from historical data. \textit{J.\ Business}, 53(1), 67--78.
\item Michaud, R.O.\ (1989). The Markowitz optimization enigma. \textit{Financial Analysts J.}, 45(1), 31--42.
\item Wolpert, D.H.\ (1992). Stacked generalization. \textit{Neural Networks}, 5(2), 241--259.
\item Higham, N.J.\ (2002). Computing the nearest correlation matrix. \textit{IMA J.\ Numer.\ Anal.}, 22(3), 329--343.
\item Ledoit, O.\ \& Wolf, M.\ (2004). A well-conditioned estimator for large-dimensional covariance matrices. \textit{J.\ Multivariate Anal.}, 88(2), 365--411.
\item Cuturi, M.\ (2013). Sinkhorn distances. \textit{Advances in NIPS}, 26.
\item Peyr\'e, G.\ \& Cuturi, M.\ (2019). Computational Optimal Transport. \textit{Found.\ Trends ML}, 11(5--6), 355--607.
\end{enumerate}

\end{document}